%% ----------------------------------------------------------------------------
%%  §4.5  Funzioni continue
%% ----------------------------------------------------------------------------
\section{Funzioni continue}
\label{sec:continue}

Come nel paragrafo precedente, consideriamo funzioni $f$ definite in un dominio $A$
costituito da un intervallo, o dall'unione finita di intervalli, e punti $x_{0}$ di $A$.
Abbiamo introdotto i limiti di funzioni soprattutto per descrivere il comportamento di
una funzione nelle vicinanze di punti «singolari», come $x_{0}=0$ per $\frac{\sen x}{x}$.
Naturalmente possiamo calcolare il limite per $x\to x_{0}$ anche se la funzione non
presenta alcuna singolarità in $x_{0}$. Ad esempio abbiamo già calcolato in
\eqref{eq:lim-sen-cos-0} i limiti
\begin{equation}\label{eq:sen-cos-in-0}
  \lim_{x\to0}\sen x=0=\sen0, \qquad \lim_{x\to0}\cos x=1=\cos0:
\end{equation}
il valore limite, per $x\to0$, è uguale al valore che si ottiene calcolando la funzione
per $x=0$. Si dice che le funzioni $\sen x$ e $\cos x$ sono \emph{continue} per $x=0$ (e
in realtà, come vedremo, sono continue in ogni $x_{0}\in\R$), in accordo con la seguente
definizione.

\begin{definizione}[Funzione continua]\label{def:continua}\index{funzione!continua}\index{continuità}
Siano $f:A\to\R$ e $x_{0}\in A$ un punto di accumulazione per $A$. La funzione $f$ si
dice \emph{continua nel punto $x_{0}$} se
\begin{equation}\label{eq:def-continua}
  \lim_{x\to x_{0}}f(x)=f(x_{0}) .
\end{equation}
La funzione si dice \emph{continua in un insieme} $B\subseteq A$ se è continua in ogni
punto $x_{0}\in B$; si dice semplicemente \emph{continua} se è continua in tutto il suo
dominio.
\end{definizione}

Nella definizione sono contenute tre richieste: che $f$ sia definita in $x_{0}$, che
esista il limite di $f$ per $x\to x_{0}$, e che il limite coincida con il valore. Se
$A=[a,b]$ e $x_{0}=a$, nella \eqref{eq:def-continua} intervengono soltanto punti $x>a$
(non ci sono punti del dominio a sinistra di $a$), e dunque il limite è di fatto il solo
limite destro $x\to a^{+}$; analogamente, se $x_{0}=b$, si considera il limite sinistro
$x\to b^{-}$. Una funzione è quindi continua nell'intervallo $[a,b]$ se è continua in ogni
punto interno, e se agli estremi il limite destro (in $a$) e il limite sinistro (in $b$)
coincidono con i valori $f(a)$ e $f(b)$.

\begin{osservazione}
Se $x_{0}$ fosse un punto isolato di $A$ (cioè un punto di $A$ che non è di
accumulazione), la \eqref{eq:def-continua} non avrebbe senso. Per convenzione si dice
che ogni funzione è continua nei punti isolati del suo dominio; con i domini che
consideriamo, intervalli non degeneri o loro unioni finite, il caso non si presenta.
\end{osservazione}

La continuità ammette due riformulazioni, che useremo continuamente: una con le
successioni, senza la restrizione $x_{n}\neq x_{0}$, e una con $\epsilon$ e $\delta$,
senza la restrizione $x\neq x_{0}$.

\begin{proposizione}[Caratterizzazioni della continuità]\label{prop:caratt-continuita}
Siano $f:A\to\R$ e $x_{0}\in A$ di accumulazione per $A$. Sono equivalenti:
\begin{enumerate}[label=\textup{(\alph*)}]
  \item $f$ è continua in $x_{0}$;
  \item per ogni successione $x_{n}$ di punti di $A$ convergente a $x_{0}$ (anche con
  $x_{n}=x_{0}$ per qualche $n$) risulta $f(x_{n})\to f(x_{0})$;
  \item per ogni $\epsilon>0$ esiste $\delta>0$ tale che $\abs{f(x)-f(x_{0})}<\epsilon$ per
  ogni $x\in A$ con $\abs{x-x_{0}}<\delta$.
\end{enumerate}
\end{proposizione}

\begin{schema}
\Ipotesi{$f:A\to\R$; \ $x_{0}\in A$ di accumulazione per $A$.}
\Tesi{(a) $\sse$ (b) $\sse$ (c).}
\end{schema}

\begin{strategia}
Si dimostra il ciclo di implicazioni (a)$\implica$(c)$\implica$(b)$\implica$(a), sfruttando
il \cref{teo:eps-delta}. L'osservazione chiave è che, quando il limite coincide con il
valore $f(x_{0})$, la condizione $x\neq x_{0}$ (o $x_{n}\neq x_{0}$) diventa superflua:
nel punto $x_{0}$ la distanza $\abs{f(x)-f(x_{0})}$ è nulla, e quindi minore di ogni
$\epsilon$.
\end{strategia}

\begin{dimostrazione}
(a)$\implica$(c): per il \cref{teo:eps-delta} la disuguaglianza
$\abs{f(x)-f(x_{0})}<\epsilon$ vale per $x\in A$ con $0<\abs{x-x_{0}}<\delta$; per
$x=x_{0}$ vale banalmente, perché il primo membro è zero.

(c)$\implica$(b): è la prima parte della dimostrazione del \cref{teo:eps-delta}, in cui
non serve più l'ipotesi $x_{n}\neq x_{0}$: se $x_{n}\to x_{0}$, definitivamente
$\abs{x_{n}-x_{0}}<\delta$, e quindi $\abs{f(x_{n})-f(x_{0})}<\epsilon$.

(b)$\implica$(a): le successioni con $x_{n}\neq x_{0}$ sono un caso particolare di quelle
considerate in (b); dunque $f(x)\to f(x_{0})$ per la \cref{def:limite-funzione}.
\end{dimostrazione}

La (b) si legge così: \emph{per una funzione continua il limite «passa dentro» la
funzione}, $\lim_{n}f(x_{n})=f\bigl(\lim_{n}x_{n}\bigr)$. La (c), scritta con
$x=x_{0}+h$, dice che $f(x_{0}+h)-f(x_{0})\to0$ per $h\to0$: a un incremento piccolo
della variabile corrisponde un incremento piccolo della funzione, che è esattamente la
definizione di Cauchy riportata nell'epigrafe del capitolo. In particolare
\begin{equation}\label{eq:continua-h}
  f \text{ è continua in } x_{0} \quad\sse\quad \lim_{h\to0}f(x_{0}+h)=f(x_{0}).
\end{equation}

\subsection{Operazioni con le funzioni continue}

\begin{teorema}[Operazioni con funzioni continue]\label{teo:operazioni-continue}
Siano $f,g:A\to\R$ e $x_{0}\in A$ di accumulazione per $A$. Se $f$ e $g$ sono continue
in $x_{0}$, sono continue in $x_{0}$ anche $f+g$, $f-g$, $f\cdot g$ e, se
$g(x_{0})\neq0$, $f/g$.
\end{teorema}

\begin{schema}
\Ipotesi{$f$, $g$ continue in $x_{0}$ \ (per il quoziente, $g(x_{0})\neq0$).}
\Tesi{$f\pm g$, $fg$, $f/g$ continue in $x_{0}$.}
\end{schema}

\begin{strategia}
Basta tradurre la continuità in termini di limiti e applicare il
\cref{teo:operazioni-funzioni} con $\ell_{1}=f(x_{0})$ e $\ell_{2}=g(x_{0})$: il limite
dell'operazione è l'operazione sui valori, cioè il valore della nuova funzione in
$x_{0}$.
\end{strategia}

\begin{dimostrazione}
Dato che il limite di somma, differenza, prodotto, quoziente è uguale rispettivamente
alla somma, differenza, prodotto, quoziente dei limiti (\cref{teo:operazioni-funzioni}),
si ha ad esempio $\lim_{x\to x_{0}}f(x)g(x)=f(x_{0})g(x_{0})$. Per il quoziente
occorre al solito fare attenzione ai punti dove il denominatore si annulla: se
$g(x_{0})\neq0$, la dimostrazione del \cref{teo:operazioni-funzioni} mostra che $g$ non
si annulla vicino a $x_{0}$, e $f(x)/g(x)\to f(x_{0})/g(x_{0})$.
\end{dimostrazione}

\begin{teorema}[Continuità della funzione composta]\label{teo:composta-continua}\index{funzione!composta!continuità}
Siano $g:X\to Y$ continua in $x_{0}\in X$ e $f:Y\to\R$ continua in $y_{0}=g(x_{0})$.
Allora la funzione composta $f\circ g$ è continua in $x_{0}$.
\end{teorema}

\begin{schema}
\Ipotesi{$g$ continua in $x_{0}$; \ $f$ continua in $g(x_{0})$.}
\Tesi{$f(g(x))\to f(g(x_{0}))$ per $x\to x_{0}$.}
\end{schema}

\begin{strategia}
Con la caratterizzazione per successioni, senza la restrizione $x_{n}\neq x_{0}$, la
dimostrazione è una catena: $x_{n}\to x_{0}$ implica $g(x_{n})\to g(x_{0})$, che a sua
volta implica $f(g(x_{n}))\to f(g(x_{0}))$. Non occorre preoccuparsi dei termini con
$g(x_{n})=g(x_{0})$, che invece creavano difficoltà nel \cref{teo:limite-composta}.
\end{strategia}

\begin{dimostrazione}
Usiamo la caratterizzazione (b) della \cref{prop:caratt-continuita}. Se $x_{n}\in X$ e
$x_{n}\to x_{0}$, per la continuità di $g$ si ha $y_{n}=g(x_{n})\to g(x_{0})=y_{0}$, con
$y_{n}\in Y$; per la continuità di $f$ in $y_{0}$ si ha allora $f(y_{n})\to f(y_{0})$,
cioè $f(g(x_{n}))\to f(g(x_{0}))$.
\end{dimostrazione}

Si noti che qui non serve l'ipotesi $g(x)\neq y_{0}$ del \cref{teo:limite-composta}:
proprio perché la (b) ammette successioni con $y_{n}=y_{0}$. Lo stesso argomento prova
la seguente versione, che useremo spesso nel calcolo dei limiti.

\begin{corollario}\label{cor:limite-composta-continua}
Siano $g:X\to Y$, $f:Y\to\R$ e $x_{0}\in\overline{\R}$ di accumulazione per $X$. Se
$\lim_{x\to x_{0}}g(x)=y_{0}$, con $y_{0}\in Y$, e $f$ è continua in $y_{0}$, allora
\begin{equation}\label{eq:limite-composta-continua}
  \lim_{x\to x_{0}}f(g(x))=f\Bigl(\lim_{x\to x_{0}}g(x)\Bigr)=f(y_{0}) .
\end{equation}
\end{corollario}

\begin{schema}
\Ipotesi{$g(x)\to y_{0}$ per $x\to x_{0}$, con $y_{0}\in Y$; \ $f$ continua in $y_{0}$.}
\Tesi{$f(g(x))\to f(y_{0})$ per $x\to x_{0}$.}
\end{schema}

\begin{strategia}
È la stessa catena del teorema precedente, con una sola differenza: ora $x_{0}$ può
non appartenere al dominio di $g$, e si considerano soltanto successioni con
$x_{n}\neq x_{0}$. Dove arrivano i valori $g(x_{n})$ non importa, perché la
caratterizzazione (b) della continuità di $f$ ammette anche $g(x_{n})=y_{0}$.
\end{strategia}

\begin{dimostrazione}
Se $x_{n}\in X\setminus\{x_{0}\}$ e $x_{n}\to x_{0}$, allora $g(x_{n})\to y_{0}$, e per
la \cref{prop:caratt-continuita}(b) applicata a $f$ (che non richiede $g(x_{n})\neq
y_{0}$) si ha $f(g(x_{n}))\to f(y_{0})$.
\end{dimostrazione}

\subsection{Continuità delle funzioni elementari}

L'importanza delle funzioni continue sta anche nel fatto che \emph{le funzioni
elementari sono continue nel loro insieme di definizione}. Esaminiamole una per volta.

\paragraph{Polinomi e funzioni razionali.} Per l'\cref{es:limite-polinomi}, ogni polinomio
è continuo su tutto $\R$, e ogni funzione razionale $P(x)/Q(x)$ è continua in ogni punto
in cui $Q(x)\neq0$.

\paragraph{Valore assoluto.} La continuità della funzione $f(x)=\abs x$ su $\R$ segue
dalla disuguaglianza triangolare inversa (\cref{cor:triangolare-inversa})
\begin{equation}\label{eq:modulo-lipschitz}
  \bigl\lvert\abs x-\abs{x_{0}}\bigr\rvert\le\abs{x-x_{0}} \qquad\perogni x,x_{0}\in\R:
\end{equation}
basta infatti porre $\delta=\epsilon$ nella \cref{prop:caratt-continuita}(c): se
$\abs{x-x_{0}}<\epsilon$, allora anche $\abs{f(x)-f(x_{0})}=\bigl\lvert\abs x-\abs{x_{0}}\bigr\rvert<\epsilon$.

\begin{proposizione}[Continuità di seno, coseno, tangente]\label{prop:sen-cos-continue}\index{funzione!seno!continuità}
Le funzioni $\sen x$ e $\cos x$ sono continue su $\R$; la funzione $\tg x$ è continua
in ogni punto $x\neq\frac\pi2+k\pi$, $k\in\Z$.
\end{proposizione}

\begin{schema}
\Ipotesi{$\sen h\to0$, $\cos h\to1$ per $h\to0$ \eqref{eq:sen-cos-in-0}; \ formule di
addizione.}
\Tesi{$\sen$, $\cos$ continue in ogni $x_{0}\in\R$; \ $\tg$ continua dove $\cos x\neq0$.}
\end{schema}

\begin{strategia}
Conosciamo la continuità di seno e coseno nel solo punto $0$ \eqref{eq:sen-cos-in-0}. Le
formule di addizione permettono di «trasportare» questa informazione in un punto
qualsiasi: $\sen(x_{0}+h)$ si esprime per mezzo di $\sen h$ e $\cos h$, che per $h\to0$
tendono a $0$ e a $1$.
\end{strategia}

\begin{dimostrazione}
Sia $x_{0}\in\R$ e sia $x_{n}\to x_{0}$. Posto $h_{n}=x_{n}-x_{0}$, si ha $h_{n}\to0$, e
quindi $\sen h_{n}\to0$, $\cos h_{n}\to1$ per il \cref{teo:lim-sen-cos}. Per la formula
di addizione del seno \eqref{eq:addizione-sen}
\begin{equation}\label{eq:sen-continuo}
\begin{split}
  \sen x_{n}=\sen(x_{0}+h_{n})&=\sen x_{0}\cos h_{n}+\cos x_{0}\sen h_{n}\\
  &\to\ \sen x_{0}\cdot1+\cos x_{0}\cdot0=\sen x_{0},
\end{split}
\end{equation}
per le operazioni con i limiti di successioni. Analogamente, per la
\eqref{eq:addizione-cos}, $\cos x_{n}=\cos x_{0}\cos h_{n}-\sen x_{0}\sen h_{n}\to\cos x_{0}$.
Per la \cref{prop:caratt-continuita}(b) seno e coseno sono continui in $x_{0}$. Infine la
continuità della funzione $\tg x=\frac{\sen x}{\cos x}$ nei punti in cui $\cos x\neq0$,
cioè per $x\neq\frac\pi2+k\pi$, discende dal \cref{teo:operazioni-continue}.
\end{dimostrazione}

\paragraph{Esponenziali, logaritmi, potenze.} Anche le funzioni esponenziali $a^{x}$
($a>0$), i logaritmi $\log_{a}x$ ($a>0$, $a\neq1$) e le potenze $x^{b}$ (per $x>0$, con
$b\in\R$) sono continue nel loro insieme di definizione. La dimostrazione richiede il
teorema sul limite delle funzioni monotòne, ed è svolta nel \cref{sec:monotone-inverse}
(\cref{teo:continuita-esponenziale}); essa non usa alcun risultato di questo paragrafo
che dipenda dalla continuità di tali funzioni, per cui possiamo servircene fin d'ora
negli esempi. Nello stesso paragrafo ricaveremo, come applicazione, alcuni limiti
notevoli che coinvolgono $e^{x}$ e $\log x$.

\begin{table}[htb]
\centering
\small
\renewcommand{\arraystretch}{1.3}
\begin{tabular}{l l l}
\toprule
funzione & continua in & dimostrazione\\
\midrule
polinomi & $\R$ & \cref{es:limite-polinomi}\\
$P(x)/Q(x)$ & $Q(x)\neq0$ & \\
$\abs x$ & $\R$ & \eqref{eq:modulo-lipschitz}\\
$\sen x$, $\cos x$ & $\R$ & \cref{prop:sen-cos-continue}\\
$\tg x$ & $x\neq\frac\pi2+k\pi$ & \\
$a^{x}$ & $\R$ & \cref{teo:continuita-esponenziale}\\
$\log_{a}x$ & $(0,+\infty)$ & \\
$x^{b}$ & $(0,+\infty)$ & \\
\bottomrule
\end{tabular}
\caption{La continuità delle funzioni elementari. Somme, prodotti, quozienti (dove il
denominatore non si annulla) e composizioni di funzioni continue sono continue
(\cref{teo:operazioni-continue,teo:composta-continua}).}
\label{tab:elementari-continue}
\end{table}

Grazie alla \cref{tab:elementari-continue} e ai \cref{teo:operazioni-continue,%
teo:composta-continua}, la continuità di funzioni come
$x\mapsto\sen(x^{2}+1)$, $x\mapsto e^{\cos x}$ o $x\mapsto\abs{\log x}$ è immediata, e i
corrispondenti limiti nei punti del dominio si calcolano per sostituzione: ad esempio
$\lim_{x\to\pi}e^{\cos x}=e^{\cos\pi}=e^{-1}$.

\begin{idea}[L'atmosfera standard: continua, ma con uno spigolo]
Nel modello dell'atmosfera standard internazionale (ISA) la temperatura, in funzione
della quota $h$ espressa in metri, è
\[
  T(h)=
  \begin{cases}
    288{,}15-0{,}0065\,h & \text{per } 0\le h\le11\,000 \ \ \text{(troposfera)},\\
    216{,}65 & \text{per } 11\,000<h\le20\,000 \ \ \text{(stratosfera)},
  \end{cases}
\]
in kelvin. Ciascun tratto è un polinomio, quindi continuo; nel punto di raccordo, la
tropopausa, il limite sinistro vale $288{,}15-0{,}0065\cdot11\,000=216{,}65$, uguale al
limite destro e al valore: $T$ è continua anche lì. Il modello non poteva essere
altrimenti: una temperatura discontinua in quota non avrebbe senso fisico, e la
pressione e la densità, che si ricavano da $T(h)$, risultano anch'esse continue. Il
grafico di $T$ ha però uno \emph{spigolo} a $11\,000$~m: la pendenza passa bruscamente da
$-6{,}5$~K/km a zero. La continuità non esclude cambi di pendenza improvvisi; li
misureremo con le derivate.
\end{idea}

\begin{notastorica}[Che cosa vuol dire «continua»?]
Per Eulero, nell'\emph{Introductio in analysin infinitorum} (1748), una funzione era
«continua» se era data da un'unica espressione analitica, e «discontinua» (o mista) se
era definita da espressioni diverse in parti diverse del dominio: secondo questo uso la
funzione $T(h)$ dell'atmosfera standard sarebbe stata discontinua, e $\frac1x$
continua! Discontinua sarebbe stata anche $\abs x$, che vale $x$ per $x\ge0$ e $-x$ per
$x<0$, ed è per noi continua ovunque. A mettere in crisi questo linguaggio fu il problema
della corda vibrante. Nel 1747 d'Alembert ne scrisse la soluzione generale come somma di
due onde che viaggiano in versi opposti; Eulero obiettò che la forma iniziale della
corda può essere qualsiasi, ad esempio triangolare, con uno spigolo, come quella di una
corda pizzicata, e che bisognava ammettere anche curve «discontinue», perfino tracciate a
mano libera. Daniel Bernoulli (1753) sostenne invece che ogni vibrazione è una
sovrapposizione di oscillazioni sinusoidali, cioè la somma di una serie trigonometrica.
Eulero e d'Alembert respinsero questa tesi: una serie di seni era, ai loro occhi, un'unica
espressione, e non poteva rappresentare una curva «discontinua». La risposta venne da
Joseph Fourier, che nella \emph{Théorie analytique de la chaleur} (1822) rappresentò con
serie trigonometriche funzioni definite a tratti, anche con salti: la stessa «formula»
poteva descrivere oggetti molto diversi, e il concetto di funzione andava separato da
quello di formula. La definizione moderna di continuità, che guarda al comportamento della
funzione vicino a ogni punto e non alla forma della formula, fu data indipendentemente da
Bernard Bolzano (1817) e da Augustin-Louis Cauchy, nel \emph{Cours d'analyse} (1821).
\end{notastorica}
